\subsection{Running Example: Moment Tensor Potentials for Ti-Al}
\label{sec:running-example}

We introduce the ontology through a concrete example: training a Moment
Tensor Potential (MTP)~\cite{shapeev2016mtp,novikov2021mlip} for the
titanium-aluminium (Ti-Al) binary alloy system. This example illustrates
how the ontology captures the full lifecycle of an MLIP study---from the
physics of the simulation cell, through the algorithm's hyperparameters
and training data, to the evaluation of the trained model.

\paragraph{The simulation problem.}
Consider a simulation cell containing $n$ atoms at positions
$\mathbf{r}_1, \ldots, \mathbf{r}_n$ with atomic types
$z_1, \ldots, z_n$ (where $z_i \in \{\text{Ti}, \text{Al}\}$ for our
Ti-Al system). The goal of an interatomic potential is to predict the
total energy $E$ and the forces $\mathbf{F}_i = -\nabla_{\mathbf{r}_i} E$
acting on each atom, given only the atomic positions and types.
Figure~\ref{fig:simulation-cell} illustrates the simulation cell and the
concept of atomic neighborhoods.

\begin{figure}[t]
\centering
\begin{subfigure}[t]{0.48\textwidth}
  \centering
  %% TODO: create 3D simulation cell figure
  \fbox{\parbox{0.9\linewidth}{\centering\vspace{3cm}%
    3D simulation cell with Ti (blue)\\and Al (orange) atoms%
    \vspace{3cm}}}
  \caption{A Ti-Al simulation cell with periodic boundary conditions.}
  \label{fig:simulation-cell-3d}
\end{subfigure}
\hfill
\begin{subfigure}[t]{0.48\textwidth}
  \centering
  %% TODO: create 2D cutoff radius diagram
  \fbox{\parbox{0.9\linewidth}{\centering\vspace{3cm}%
    2D cross-section showing $R_{\text{cut}}$\\and neighborhood $\neigh{i}$%
    \vspace{3cm}}}
  \caption{The neighborhood $\neigh{i}$ of atom~$i$: all atoms within
    $R_{\text{cut}}$ contribute to $V(\neigh{i})$. Atoms beyond
    $R_{\text{cut}}$ (grayed out) are excluded.}
  \label{fig:simulation-cell-cutoff}
\end{subfigure}
\caption{(a)~A Ti-Al simulation cell. (b)~The cutoff radius
  $R_{\text{cut}}$ defines the neighborhood $\neigh{i}$, which
  includes relative positions $\mathbf{r}_{ij}$ and atomic types $z_j$
  of the neighbors.}
\label{fig:simulation-cell}
\end{figure}

In quantum mechanics, the total energy can be computed exactly (within
the approximations of density functional theory) by solving the
Kohn--Sham equations~\cite{kohn1965dft}. However, this is
computationally expensive---scaling as $O(n^3)$ with the number of
atoms---limiting DFT to systems of a few hundred atoms. Machine learning
interatomic potentials approximate the DFT energy surface at a fraction
of the cost, enabling simulations of millions of atoms.

\paragraph{The MTP algorithm.}
In the MTP framework~\cite{shapeev2016mtp}, the total energy of an
atomic configuration $\text{cfg}$ is decomposed into local contributions:
%
\begin{equation}\label{eq:mtp-energy}
  E^{\text{mtp}}(\text{cfg}) = \sum_{i=1}^{n} V(\neigh{i}),
\end{equation}
%
where $\neigh{i}$ is the \emph{neighborhood} of atom~$i$, defined as
the set of all atoms within the cutoff radius $R_{\text{cut}}$. The
neighborhood consists of the atomic type $z_i$ of the central atom, and
for each neighbor $j$ within $R_{\text{cut}}$: the atomic type $z_j$ and
the relative position vector $\mathbf{r}_{ij} = \mathbf{r}_j -
\mathbf{r}_i$. Crucially, atoms beyond $R_{\text{cut}}$ do not
contribute to $V(\neigh{i})$---this \emph{locality assumption} is
what makes the potential computationally efficient.

The local energy $V(\neigh{i})$ is expressed as a linear combination
of basis functions:
%
\begin{equation}\label{eq:mtp-basis}
  V(\neigh{i}) = \sum_{\alpha} \xi_\alpha \, B_\alpha(\neigh{i}),
\end{equation}
%
where $B_\alpha$ are the \emph{moment tensor descriptors}---invariant
functions of the neighborhood that capture the local atomic
environment---and $\boldsymbol{\xi} = \{\xi_\alpha\}$ are the
\emph{parameters to be learned} by fitting to DFT reference data. The
basis functions $B_\alpha$ are constructed from contractions of moment
tensors $M_{\mu,\nu}$ involving radial functions $Q^{(\beta)}(|\mathbf{r}_{ij}|)$
that vanish smoothly at $R_{\text{cut}}$.

\paragraph{Hyperparameters.}
The MTP algorithm has several hyperparameters that control the model's
expressiveness and computational cost:
%
\begin{itemize}
\item $R_{\text{cut}}$ (\emph{cutoff radius}): determines the size of
  atomic neighborhoods. A larger cutoff captures longer-range
  interactions but increases computational cost.
\item $R_{\text{min}}$ (\emph{minimum distance}): the shortest expected
  interatomic distance, used to define the radial basis domain.
\item $N_Q$ (\emph{number of radial basis functions}): controls the
  resolution of the radial description.
\item $\text{lev}_{\text{max}}$ (\emph{maximum level}): determines which
  basis functions $B_\alpha$ are included, controlling the trade-off
  between accuracy and model size.
\end{itemize}
%
These are \emph{not} learned from data---they are chosen by the
researcher before training. In contrast, the coefficients
$\boldsymbol{\xi}$ are determined by fitting
Equation~\eqref{eq:mtp-basis} to DFT-computed energies, forces, and
stresses via linear regression.

\paragraph{Training data.}
The ground truth for training comes from DFT calculations: for a set of
atomic configurations of the Ti-Al system, one computes the total energy,
forces on each atom, and (optionally) stresses. Each configuration is a
snapshot of atomic positions in a periodic cell, generated by sampling
from molecular dynamics trajectories, random perturbations, or active
learning~\cite{novikov2021mlip}. A typical training dataset for Ti-Al
might contain 15{,}000 configurations computed with the PBE
exchange-correlation functional, PAW pseudopotentials, and a plane-wave
energy cutoff of 520~eV using the VASP code.

\paragraph{Ontology mapping.}
This example maps directly onto the \MLIPsOntology{} as follows.
The MTP algorithm is an instance of $\MLIPAlgorithmC$ with four
$\HyperparameterC$ instances ($R_{\text{cut}}$, $R_{\text{min}}$, $N_Q$,
$\text{lev}_{\text{max}}$). The Ti-Al material is a $\MaterialSystemC$.
The 15{,}000 configurations form a $\TrainingDatasetC$, linked to a
$\DFTCalculationC$ with $\DFTSettingsC$ recording the PBE functional and
PAW pseudopotentials.

A $\TrainingRunC$ records the concrete training event: it $\executesR$
the MTP algorithm, $\runsOnR$ the Ti-Al dataset, and $\producesR$ a
$\TrainedModelC$---the specific set of coefficients $\boldsymbol{\xi}$
fitted to this data with $R_{\text{cut}} = 5.0$~\AA{},
$\text{lev}_{\text{max}} = 16$, and $N_Q = 8$. These concrete
values are recorded as $\HyperparameterSettingC$ instances.

When the trained model is evaluated on a test set (e.g., computing the
RMSE of predicted energies against DFT reference values), the
evaluation is captured as a $\BenchmarkResultC$ that $\evaluatesModelR$
the trained model, $\targetMaterialR$ the Ti-Al system, and reports
$\AccuracyMetricC$ instances (e.g., RMSE of energy = 1.2~meV/atom).

Figure~\ref{fig:mtp-ontology-mapping} shows how these ontology instances
relate to each other through the ontology's roles.

\begin{figure}[t]
\centering
%% TODO: create instance diagram showing the MTP/Ti-Al example
\fbox{\parbox{0.9\textwidth}{\centering\vspace{3cm}%
  Instance diagram: MTP algorithm → training run → trained model → benchmark result,\\
  with hyperparameter settings, DFT calculation, and accuracy metrics.%
  \vspace{3cm}}}
\caption{Ontology instance diagram for the running example: training an
  MTP model for Ti-Al. Yellow nodes are instances; blue nodes are classes
  from the \MLIPsOntology{}. The diagram shows how the MTP algorithm,
  training data, trained model, and benchmark results are connected
  through the ontology's roles.}
\label{fig:mtp-ontology-mapping}
\end{figure}
